% Fragmento público generado desde propietarios íntegros.
% Residencia: 52.7.2.
% Fuente: ../incoming/radion_monografia_20260827/manuscrito/sections/08_complejo_electromagnetismo.tex:59-158
\subsubsection{Opérateur constitutif sur les canaux \(90/120\)}

\paragraph{Décomposition en composantes commune et orientée}

On définit
\begin{align}
\mathcal E
&=\log\bigl[(1-q_+^{90})(1-q_-^{90})\bigr]
-\log\bigl[(1-q_+^{120})(1-q_-^{120})\bigr],
\label{eq:Ecommon}\\
\mathcal B
&=\log\frac{1-q_-^{90}}{1-q_+^{90}}
-\log\frac{1-q_-^{120}}{1-q_+^{120}}.
\label{eq:Borient}
\end{align}
Sous \(C^*\mapsto-C^*\), \(q_+\leftrightarrow q_-\) s'échangent. Par
conséquent, \(\mathcal E\) est pair et \(\mathcal B\) est impair.

Les lois constitutives normalisées sont
\begin{equation}
\widehat\mu=e^{\mathcal E+\mathcal B},
\qquad
\widehat\varepsilon=e^{\mathcal E-\mathcal B},
\qquad
\widehat Z=e^{\mathcal B},
\qquad
\widehat c=e^{-\mathcal E}.
\label{eq:constitutives}
\end{equation}
On vérifie exactement
\begin{equation}
\widehat Z=\sqrt{\frac{\widehat\mu}{\widehat\varepsilon}},
\qquad
\widehat c=(\widehat\mu\widehat\varepsilon)^{-1/2}.
\label{eq:constitutive-identities}
\end{equation}
La conjugaison produit
\begin{equation}
C^*\mapsto-C^*:\qquad
\widehat\mu\leftrightarrow\widehat\varepsilon,
\quad
\widehat Z\leftrightarrow\widehat Z^{-1},
\quad
\widehat c\mapsto\widehat c.
\label{eq:EM-conjugation}
\end{equation}

\begin{theorem}[Polarité constitutive]
Le lecteur \(90/120\) transforme le mode commun et le contraste orienté du
couple \((A,C^*)\) en deux lois constitutives conjuguées. L'inversion de feuillet
échange perméabilité et permittivité, inverse l'impédance et conserve la
vitesse normalisée.
\end{theorem}

Telle est la formulation précise de l'intuition \enquote{partie électrique et
partie magnétique}. On n'écrit ni \(A=E\) ni \(C^*=B\). On écrit l'application
\[
(A,C^*)\to(x,y)\to(q_+,q_-)
\to(\mathcal E,\mathcal B)
\to(\widehat\varepsilon,\widehat\mu,\widehat Z,\widehat c).
\]

\paragraph{Réalisation élastique--rotationnelle de la polarité constitutive}

La même information peut être publiée dans le régime complexe au moyen de
\begin{equation}
Z_{\mathrm{em}}=e^{\mathcal E/2}
\left(\cos\frac{\mathcal B}{2}\,I
-\sin\frac{\mathcal B}{2}\,J\right).
\label{eq:Zem}
\end{equation}
Le facteur \(e^{\mathcal E/2}\) est une dilatation élastique ; le facteur généré
par \(J\) est une rotation. L'écriture complexe n'ajoute pas une dimension
imaginaire mystérieuse à un monde réel préalablement donné : elle enregistre deux modes
opératoires que le TRIT a déjà séparés.

\paragraph{Transduction de l'orientation en charge effective}

Le feuillet \(\sigma\) du radion fixe une orientation avant d'introduire une
unité de charge. La mécanique dimensionnelle dispose ensuite le transducteur
\[
e_{\mathrm{int}}=\sqrt{\frac{2\alphah\hfull}{Z_{\mathrm{int}}}},
\qquad
Q=n_Qe_{\mathrm{int}}.
\]
La chaîne causale est
\[
\text{feuillet}\to\text{orientation}\to\text{constitutive}
\to\text{transducteur dimensionnel}\to\text{charge physique}.
\]
Appeler l'unité \emph{radion} n'identifie pas prématurément l'orientation
au coulomb ; cela rend visible le lieu où le signe de charge peut résider.

\begin{narrativebox}
L'électricité et le magnétisme ne sont pas collés au cercle depuis l'extérieur. Ils émergent
lorsque le même état est lu à travers le produit et le quotient : le produit
retient le mode commun, le quotient retient l'orientation. L'oscillateur LC
réalise l'échange ; le lecteur \(90/120\) publie les lois constitutives ; la
carte dimensionnelle attribue des unités.
\end{narrativebox}
